\begin{scanpage}{0062}{45}
\[\left.\begin{aligned}E_\theta&=H_\phi=ikA_\theta,\\E_\phi&=-H_\theta=ikA_\phi,\\E_R&=H_R=0,\end{aligned}\right\}\tag{7.08}\]
where
\[\left.\begin{aligned}A_\phi&=A_y\cos\phi-A_x\sin\phi,\\A_\theta&=(A_x\cos\phi+A_y\sin\phi)\cos\theta-A_z\sin\theta.\end{aligned}\right\}\tag{7.09}\]
Substituting here the values
\[\left.\begin{aligned}A_y&=-H_{0x}\frac a{k\sin\theta}J_1(ka\sin\theta)\frac{e^{ikR}}R,\\A_x&=A_z=0,\end{aligned}\right\}\tag{7.10}\]
which result from Equations (7.02) and (7.06), let us find the field radiated by the uniform part of the current in the form
\[\left.\begin{aligned}E_\theta=H_\phi&=-iaH_{0x}\frac{\sin\phi}{\sin\theta}\cos\theta J_1(ka\sin\theta)\frac{e^{ikR}}R,\\E_\phi=-H_\theta&=-iaH_{0x}\frac{\cos\phi}{\sin\theta}J_1(ka\sin\theta)\frac{e^{ikR}}R.\end{aligned}\right\}\tag{7.11}\]
The function $J_1(ka\sin\theta)$ is a first order Bessel function. By using its asymptotic expression
\[J_1(ka\sin\theta)=\sqrt{\frac2{\pi ka\sin\theta}}\cos\left(ka\sin\theta-\frac{3\pi}4\right),\tag{7.12}\]
which is applicable when $ka\sin\theta\gg1$, one is able to rewrite Equations (7.11) in the following form:
\[\left.\begin{aligned}E_\theta=H_\phi&=-iH_{0x}\sqrt{\frac a{2\pi k\sin\theta}}\cos\theta\frac{\sin\phi}{\sin\theta}\times\\&\quad\times\left[e^{-i(ka\sin\theta-\frac{3\pi}4)}+e^{i(ka\sin\theta-\frac{3\pi}4)}\right]\frac{e^{ikR}}R,\\E_\phi=-H_\theta&=-iH_{0x}\sqrt{\frac a{2\pi k\sin\theta}}\frac{\cos\phi}{\sin\theta}\times\\&\quad\times\left[e^{-i(ka\sin\theta-\frac{3\pi}4)}+e^{i(ka\sin\theta-\frac{3\pi}4)}\right]\frac{e^{ikR}}R.\end{aligned}\right\}\tag{7.13}\]
\end{scanpage}
